\subsection{Worked example: bounded eventuality with an exclusion window}
Let
\[
  f=(\Eventually_{\le5}p)\land(\Always_{<2}\neg p).
\]
It requires $p$ to be false at the origin and until age two, then true by age
five. Fig.~\ref{fig:translation-reference-tba} shows a hand-simplified
equivalent one-clock TBA, not the syntax-directed output, which assigns
separate clocks to the two distinct timed subformulas. The steps below show
the clocked-LTL translation and timed reinterpretation. This is a construction
illustration: the saved Spot graph and its two-clock TBA transcription are not
covered by the proved Rocq backend contract. The concrete instance
checked in Rocq is the alarm-response example in
Section~\ref{sec:alarm-example}.

\begin{figure}[t]
\centering
\begin{tikzpicture}[node distance=22mm]
  \node[fmstate] (q0) {$q_{\mathit{init}}$};
  \node[fmstate,right=of q0] (q1) {$q_{\mathit{wait}}$};
  \node[fmaccept,right=of q1] (q2) {$q_{\mathit{acc}}$};
  \draw[fmedge] ([xshift=-8mm]q0.west) -- (q0.west);
  \draw[fmedge] (q0) -- node[fmlabel,above] {$\neg p\;/\;\mathsf{rst}(x)$} (q1);
  \draw[fmedge] (q1) edge[loop above] node[fmlabel] {$x\le5,\;\neg p$} (q1);
  \draw[fmedge] (q1) -- node[fmlabel,above] {$2\le x\le5,\;p$} (q2);
  \draw[fmedge] (q2) edge[loop above] node[fmlabel] {$\top$} (q2);
\end{tikzpicture}

\caption{Hand-simplified equivalent event-guarded TBA for $f$. The initial
transition requires $\neg p$ and resets $x$; in the waiting location,
$x\le5$, and acceptance also requires $x\ge2$. Thus $p$ first holds at age
in $[2,5]$.}
\label{fig:translation-reference-tba}
\end{figure}

\begin{enumerate}[label=Step~\arabic*:,leftmargin=*]
\item The formula is in negation normal form. Expanding the derived operators
and applying the timed-operator identities gives
\[
\begin{aligned}
f
 &= (\top\Until_{\le5}p)\land(\bot\Release_{<2}\neg p)\\
 &= \bigl(p\lor(\top\land\top\hUntil_{\le5}p)\bigr)\\
 &\qquad\land\bigl(\neg p\land(\bot\lor\bot\hRelease_{<2}\neg p)\bigr).
\end{aligned}
\]
These current-position cases follow Eqs.~\eqref{eq:derive-ule} and
\eqref{eq:derive-rle}.

\item The two distinct primitive timed subformulas receive clocks $x$ and
$y$, respectively. Their clocked-LTL translation is
\begin{align*}
\Phi_U={}&p\lor\bigl(\top\land\Next\Bigl(
  ((x\le5)\land\Unch(x)\land\top)\Until\\[-1mm]
&\qquad((x\le5)\land p)\Bigr)\bigr),\\[-1mm]
\Phi_R={}&\neg p\land\bigl(\bot\lor(\Reset(y)\land\\[-1mm]
&\qquad\Next(\neg p\WeakUntil((y\ge2)\lor(\bot\land\neg p))))\bigr),\\[-1mm]
\T(f)={}&\Phi_U\land\Phi_R.
\end{align*}
The Until clause preserves $x$; the Release test $y\ge2$ complements the source
bound $<2$, allowing $p$ from age two.

\item The saved Spot 2.16 graph for $\T(f)$ is redrawn on the left of
Fig.~\ref{fig:translation-steps45} (top). The clock guards $x\le5$ and
$y\ge2$, the markers $\mathsf{unch}(x)$ and $\mathsf{rst}(y)$, and the action
atom $p$ are propositions to the backend, and the negated action atom
$\neg p=\mathsf{LNAct}(p)$ reaches Spot as the negation of the atom $p$
(the default of the tool; Section~\ref{sec:tba-construction}). No cube of
the graph therefore contains both $p$ and $\neg p$, and every edge is
realizable by a timed event. The graph has four
states and seven transitions, the size reported for F11 in
Table~\ref{tab:compiled}.

\item Relaxation removes negative literals over extended clock atoms, here
$\neg(y\ge2)$. Reset completion resets the preservation clock $x$ unless
$\mathsf{unch}(x)$ is present; the restart clock $y$ is reset exactly when
$\mathsf{rst}(y)$ is present. The bottom of
Fig.~\ref{fig:translation-steps45} shows the same state graph after these
operations: each edge carries the event, the clock guard, and the resets. The
event is an action ($p$ or $\neg p$), and $\top$ stands for any event.
\end{enumerate}

\begin{figure*}[!t]
\centering
\begin{tikzpicture}[x=1cm,y=1cm]
% Left: the saved Spot graph for T(f) (negated event atom printed as !"p").
% Right of the arrows: guards and resets of the reset-completed TBA.
\begin{scope}[xshift=0cm]
  \node[fmnote,font=\small\bfseries,anchor=west] at (-0.9,1.55) {Spot output};
  \node[fmstate] (a0) at (0.0,0.0) {$1$};
  \node[fmstate] (a1) at (4.6,0.0) {$2$};
  \node[fmstate] (a3) at (9.8,0.0) {$3$};
  \node[fmaccept] (a2) at (4.6,-2.7) {$0$};
  \draw[fmedge] (-1.0,0.0) -- (a0);
  \draw[fmedge] (a0) -- node[fmlabel,above] {$\neg p\land\mathsf{rst}(y)$} (a1);
  \draw[fmedge] (a1) edge[loop above,looseness=6]
    node[fmlabel] {$\neg p\land\mathsf{unch}(x)\land x\le5\land\neg(y\ge2)$} (a1);
  \draw[fmedge] (a1) -- node[fmlabel,below,yshift=-1pt] {$\neg p\land\mathsf{unch}(x)\land x\le5\land y\ge2$} (a3);
  \draw[fmedge] (a1) -- node[fmlabel,left] {$p\land x\le5\land y\ge2$} (a2);
  \draw[fmedge] (a3) edge[loop above,looseness=6]
    node[fmlabel] {$\neg p\land\mathsf{unch}(x)\land x\le5$} (a3);
  \draw[fmedge] (a3) -- node[fmlabel,above right,xshift=-2pt] {$p\land x\le5$} (a2);
  \draw[fmedge] (a2) edge[loop below,looseness=5] node[fmlabel] {$\top$} (a2);
\end{scope}

\begin{scope}[yshift=-6.6cm]
  \node[fmnote,font=\small\bfseries,anchor=west] at (-0.9,1.55) {Reset-completed TBA};
  \node[fmstate] (b0) at (0.0,0.0) {$1$};
  \node[fmstate] (b1) at (4.6,0.0) {$2$};
  \node[fmstate] (b3) at (9.8,0.0) {$3$};
  \node[fmaccept] (b2) at (4.6,-2.7) {$0$};
  \draw[fmedge] (-1.0,0.0) -- (b0);
  \draw[fmedge] (b0) -- node[fmlabel,above] {$\neg p\;/\;\mathsf{rst}(x,y)$} (b1);
  \draw[fmedge] (b1) edge[loop above,looseness=6]
    node[fmlabel] {$\neg p,\;x\le5$} (b1);
  \draw[fmedge] (b1) -- node[fmlabel,below,yshift=-1pt] {$\neg p,\;x\le5\land y\ge2$} (b3);
  \draw[fmedge] (b1) -- node[fmlabel,left] {$p,\;x\le5\land y\ge2\;/\;\mathsf{rst}(x)$} (b2);
  \draw[fmedge] (b3) edge[loop above,looseness=6]
    node[fmlabel] {$\neg p,\;x\le5$} (b3);
  \draw[fmedge] (b3) -- node[fmlabel,above right,xshift=-2pt] {$p,\;x\le5\;/\;\mathsf{rst}(x)$} (b2);
  \draw[fmedge] (b2) edge[loop below,looseness=5] node[fmlabel] {$\top\;/\;\mathsf{rst}(x)$} (b2);
\end{scope}
\end{tikzpicture}

\caption{The saved Spot 2.16 graph for $\T(f)$ (top, redrawn; the labels are
propositional cubes over the action atom $p$, the clock guards, and the markers
$\mathsf{unch}(x)$ and $\mathsf{rst}(y)$) and the reset-completed TBA obtained
from it by hand (bottom; same state numbers). A TBA edge is labeled
by the event, the clock guard, and, after a slash, the clocks reset; $\top$ is
any event, and an edge without a slash resets no clock. This graph-to-TBA
illustration is not the formally verified alarm-response backend instance.}
\label{fig:translation-steps45}
\end{figure*}
